Entre 1960 y 1970, un grupo de \textit{hackers}\footnote{El término ``hácker'' ha sido tergiversado a través de los tiempos. En este artículo, se usará el término ``hácker'', refiriéndose al espíritu original, tal como nació en el MIT, durante 1961. Es decir, refiriéndose a aquellas personas que disfrutan de jugar con la tecnología, amoldándola y, en ocasiones, rompiéndola para que funcione de una forma en la que no fue planeada.}, leyendo publicaciones sobre el funcionamiento de las frecuencias usadas para realizar llamadas hechas por Bell, descubrieron cómo funcionaban los teléfonos. \cite{bluebox:sarts}

Con insasiable sed de conocimiento, y con bastos conocimientos en electrónica adquiridos durante muchas jornadas de \textit{hacking}, crean un aparato que permite realizar llamadas de forma gratuita: \textit{La Caja Azul}.

\begin{figure}
\centering
\epsfig{file=assets/bluebox.eps}
\caption{Un \textit{Blue Box} en el museo de \textit{Powerhouse}.
Autor: Maksym Kozienko}
\end{figure}

Varios meses después, Bell se da cuenta de las consecuencias que toda esta información publicada podrían causar, y destruyó físicamente las páginas en las que se encontraba toda esta información, sin embargo ya era muy tarde, porque los háckers ya habían digerido el conocimiento.

Indudablemente, los dueños de las empresas de telecomunicaciones estaban enojados y buscaban culpables, y descargaron toda su ira en los pocos háckers que pudieron atrapar e incriminar. Entre ellos estaba Kevin Mitnick, cuya historia es muy conocida. Kevin Mitnick fue acusado de ser capaz de desencadenar una guerra nuclear si se le permitía acceso a un teléfono.

La ira de estos empresarios estaba parcialmente justificada. Darle mantenimiento a todo este sistema telefónico era caro e ineficaz por entonces. El internet a la velocidad que hoy se consideraría ``normal'', por entonces era un lujo al que solo personas con el conocimiento suficiente, y con el dinero suficiente podían acceder.\cite{bluebox:documentary}

Hoy en día, con el internet a la vuelta de la esquina, el beneficio de las ``llamadas gratuitas'' ya no perjudica a nadie, dando a luz al VoIP: \textit{Voice over Internet Protocol}.

Pese a esto, muchas compañías de telefonía aún ven a las ``llamadas gratuitas'' como una abominable blasfemia a su seguridad financiera, y cuál si fuera el fruto prohibido del Edén, intentan restringir y regular su uso.
